% p327_1; source PDF page 327; manually verified editable table.
\begin{tikzpicture}[x=1cm,y=-1cm,inner sep=0pt,outer sep=0pt]
\path[use as bounding box] (0,0) rectangle (16.000,20.920);
\draw[line width=.4pt] (0.000,0.000) rectangle (3.850,0.815);
\node[align=center,text width=3.590cm,font={\linespread{1}\fontsize{10}{13.50}\selectfont\bfseries}] at (1.925,0.407) {罪数关系};
\draw[line width=.4pt] (3.850,0.000) rectangle (7.700,0.815);
\node[align=center,text width=3.590cm,font={\linespread{1}\fontsize{10}{13.50}\selectfont\bfseries}] at (5.775,0.407) {行为类型};
\draw[line width=.4pt] (7.700,0.000) rectangle (16.000,0.815);
\node[align=center,text width=8.040cm,font={\linespread{1}\fontsize{10}{13.50}\selectfont\bfseries}] at (11.850,0.407) {范例};
\draw[line width=.4pt] (0.000,0.815) rectangle (3.850,8.045);
\node[align=center,text width=3.590cm,font={\linespread{1}\fontsize{10}{13.50}\selectfont}] at (1.925,4.430) {危险驾驶罪与交通肇事罪};
\draw[line width=.4pt] (3.850,0.815) rectangle (7.700,4.030);
\node[align=center,text width=3.590cm,font={\linespread{1}\fontsize{10}{13.50}\selectfont}] at (5.775,2.422) {一个危险流};
\draw[line width=.4pt] (7.700,0.815) rectangle (16.000,4.030);
\node[align=left,text width=8.040cm,font={\linespread{1}\fontsize{10}{13.50}\selectfont}] at (11.850,2.422) {甲醉驾行驶了10公里，在第11公里因醉驾而过失撞死行人。第一阶段，甲已构成危险驾驶罪。第二阶段，甲又构成交通肇事罪。由于是同一个危险流（醉驾制造的危险流），重罪吸收轻罪，最终只定交通肇事罪};
\draw[line width=.4pt] (3.850,4.030) rectangle (7.700,8.045);
\node[align=center,text width=3.590cm,font={\linespread{1}\fontsize{10}{13.50}\selectfont}] at (5.775,6.037) {两个危险流};
\draw[line width=.4pt] (7.700,4.030) rectangle (16.000,8.045);
\node[align=left,text width=8.040cm,font={\linespread{1}\fontsize{10}{13.50}\selectfont}] at (11.850,6.037) {甲醉驾行驶了10公里，在第11公里，不是因为醉驾，而是因为新的违章行为（如闯红灯）过失撞死行人。第一阶段，甲已构成危险驾驶罪。第二阶段，甲又构成交通肇事罪。由于是两个不同的危险流（醉驾制造的危险流和闯红灯制造的危险流），应数罪并罚};
\draw[line width=.4pt] (0.000,8.045) rectangle (3.850,12.875);
\node[align=center,text width=3.590cm,font={\linespread{1}\fontsize{10}{13.50}\selectfont}] at (1.925,10.460) {危险驾驶罪与以危险方法危害公共安全罪};
\draw[line width=.4pt] (3.850,8.045) rectangle (7.700,10.060);
\node[align=center,text width=3.590cm,font={\linespread{1}\fontsize{10}{13.50}\selectfont}] at (5.775,9.052) {一个行为};
\draw[line width=.4pt] (7.700,8.045) rectangle (16.000,10.060);
\node[align=left,text width=8.040cm,font={\linespread{1}\fontsize{10}{13.50}\selectfont}] at (11.850,9.052) {甲醉酒在高速公路上逆向行驶。甲的一个行为同时触犯危险驾驶罪和以危险方法危害公共安全罪，想象竞合，择一重罪论处，定以危险方法危害公共安全罪};
\draw[line width=.4pt] (3.850,10.060) rectangle (7.700,12.875);
\node[align=center,text width=3.590cm,font={\linespread{1}\fontsize{10}{13.50}\selectfont}] at (5.775,11.467) {两个行为};
\draw[line width=.4pt] (7.700,10.060) rectangle (16.000,12.875);
\node[align=left,text width=8.040cm,font={\linespread{1}\fontsize{10}{13.50}\selectfont}] at (11.850,11.467) {甲醉驾行驶了5公里，被警察发现，为了逃跑，横冲直撞，撞死行人。甲的第一个行为构成危险驾驶罪，第二个行为构成以危险方法危害公共安全罪，数罪并罚};
\draw[line width=.4pt] (0.000,12.875) rectangle (3.850,18.205);
\node[align=center,text width=3.590cm,font={\linespread{1}\fontsize{10}{13.50}\selectfont}] at (1.925,15.540) {交通肇事罪与以危险方法危害公共安全罪};
\draw[line width=.4pt] (3.850,12.875) rectangle (7.700,15.290);
\node[align=left,text width=3.590cm,font={\linespread{1}\fontsize{10}{13.50}\selectfont}] at (5.775,14.082) {先构成交通肇事罪，后构成以危险方法危害公共安全罪};
\draw[line width=.4pt] (7.700,12.875) rectangle (16.000,15.290);
\node[align=left,text width=8.040cm,font={\linespread{1}\fontsize{10}{13.50}\selectfont}] at (11.850,14.082) {并罚。例如，甲先构成交通肇事罪，后逃逸中横冲直撞，构成以危险方法危害公共安全罪。两罪并罚};
\draw[line width=.4pt] (3.850,15.290) rectangle (7.700,18.205);
\node[align=left,text width=3.590cm,font={\linespread{1}\fontsize{10}{13.50}\selectfont}] at (5.775,16.747) {先构成以危险方法危害公共安全罪，后构成交通肇事罪};
\draw[line width=.4pt] (7.700,15.290) rectangle (16.000,18.205);
\node[align=left,text width=8.040cm,font={\linespread{1}\fontsize{10}{13.50}\selectfont}] at (11.850,16.747) {吸收。例如，甲驾车在高速公路上高速逆行，10公里后不慎撞上汽车，导致对方死亡。这是一个危险流演变为实害结果。以危险方法危害公共安全罪吸收交通肇事罪};
\draw[line width=.4pt] (0.000,18.205) rectangle (3.850,20.920);
\node[align=center,text width=3.590cm,font={\linespread{1}\fontsize{10}{13.50}\selectfont}] at (1.925,19.562) {危险驾驶罪与运输类犯罪（运输假币罪，运输毒品罪）};
\draw[line width=.4pt] (3.850,18.205) rectangle (7.700,20.920);
\node[align=center,text width=3.590cm,font={\linespread{1}\fontsize{10}{13.50}\selectfont}] at (5.775,19.562) {一个行为};
\draw[line width=.4pt] (7.700,18.205) rectangle (16.000,20.920);
\node[align=left,text width=8.040cm,font={\linespread{1}\fontsize{10}{13.50}\selectfont}] at (11.850,19.562) {甲醉酒驾车运输毒品，一个驾驶运输行为同时触犯危险驾驶罪和运输毒品罪，想象竞合，择一重罪论处，定运输毒品罪};
\end{tikzpicture}
